\documentclass[11pt]{article}

\usepackage[margin=1in]{geometry}
\usepackage{amsmath, amssymb}
\usepackage{xcolor}
\usepackage{listings}
\usepackage{hyperref}
\usepackage{array}
\usepackage{booktabs}

\hypersetup{
    colorlinks=true,
    linkcolor=blue!60!black,
    urlcolor=blue!60!black,
    citecolor=blue!60!black
}

\lstdefinestyle{bugpython}{
  language=Python,
  basicstyle=\ttfamily\small,
  keywordstyle=\color{blue!60!black},
  stringstyle=\color{green!40!black},
  commentstyle=\color{gray!70!black},
  showstringspaces=false,
  breaklines=true,
  frame=single,
  rulecolor=\color{gray!30},
  columns=fullflexible,
  keepspaces=true
}

\title{SymPy Bug Report}
\date{}

\begin{document}

\author{}
\maketitle

\section{Bug 50: Wrong negative-branch antiderivative for \texorpdfstring{$\operatorname{asec}(x)$}{asec(x)}}

\subsection{Minimal Reproducer}

\medskip

\begin{lstlisting}[style=bugpython]
from sympy import N, asec, diff, integrate, simplify, symbols

x = symbols("x", real=True)
residual = diff(integrate(asec(x), x), x) - asec(x)

print("residual at x=-2:", simplify(residual.subs(x, -2)))
print("numeric residual:", N(residual.subs(x, -2), 50))
\end{lstlisting}

\subsection{Observed Behavior}

\medskip

\begin{lstlisting}[style=bugpython]
residual at x=-2: -2*sqrt(3)/3
numeric residual: -1.1547005383792515290182975610039149112952035025403
\end{lstlisting}

The expression returned by \texttt{integrate(asec(x), x)} therefore does not
differentiate back to \texttt{asec(x)} at the nonsingular real point \(x=-2\).
The residual is exact and nonzero, so this is not merely a presentation issue
or a failure to simplify to a different but equivalent form.

\subsection{Expected Behavior}

For real \(x\) on the smooth interval \(x<-1\), any antiderivative returned for
\(\operatorname{asec}(x)\) should satisfy
\[
  \frac{d}{dx}F(x)=\operatorname{asec}(x).
\]
In particular,
\[
  \left.
  \left(\frac{d}{dx}\operatorname{integrate}(\operatorname{asec}(x),x)
  -\operatorname{asec}(x)\right)\right|_{x=-2}
  = 0.
\]

\subsection{Mathematical Explanation}

Using integration by parts gives
\[
  \int \operatorname{asec}(x)\,dx
  = x\operatorname{asec}(x)
    - \int \frac{1}{x\sqrt{1-1/x^2}}\,dx .
\]
On the negative real branch \(x<-1\),
\[
  \sqrt{1-\frac{1}{x^2}}=\frac{\sqrt{x^2-1}}{|x|}
  =-\frac{\sqrt{x^2-1}}{x},
\]
so the remaining radical integrand is
\[
  \frac{1}{x\sqrt{1-1/x^2}}
  =-\frac{1}{\sqrt{x^2-1}}.
\]
A real primitive for this subintegrand on \(x<-1\) is
\(\operatorname{acosh}(|x|)\), since
\[
  \frac{d}{dx}\operatorname{acosh}(|x|)
  =-\frac{1}{\sqrt{x^2-1}}.
\]
Thus a branch-correct antiderivative on this interval is
\[
  x\operatorname{asec}(x)-\operatorname{acosh}(|x|)+C.
\]

The diagnosed SymPy expression instead contains the real factor
\[
  \frac{x}{|x|}\operatorname{acosh}(|x|),
\]
whose derivative has the opposite sign on \(x<-1\).  Subtracting that wrong
subprimitive from the integration-by-parts result makes the final derivative
\[
  \operatorname{asec}(x)-\frac{2}{\sqrt{x^2-1}}
\]
on the negative branch.  At \(x=-2\), this gives the observed residual
\(-2/\sqrt{3}=-2\sqrt{3}/3\).

\subsection{Source-Level Diagnosis}

The diagnosis identifies the branch error in the Meijer-G and hyperexpand path,
not in \texttt{asec.fdiff} or in the integration-by-parts rule itself.  The
public call \texttt{integrate(asec(x), x)} first uses manual integration by
parts, producing
\[
  x\operatorname{asec}(x)
  - \int \frac{1}{x\sqrt{1-1/x^2}}\,dx .
\]
The default integration machinery in
\nolinkurl{sympy/integrals/integrals.py} then recursively evaluates the
remaining unevaluated subintegral.  Risch and heurisch do not evaluate that
radical integral, so the Meijer path calls \texttt{meijerint\_indefinite};
\texttt{\_rewrite1} rewrites the radical integrand as a Meijer-G term, and
\texttt{hyperexpand} expands the integrated term.

The source-level notes locate the problematic continuation in
\nolinkurl{sympy/functions/special/hyper.py}.  In particular,
\texttt{HyperRep.\_eval\_rewrite\_as\_nonrep} chooses a split based only on
\(\lvert x\rvert>1\), while \texttt{HyperRep\_asin1.\_expr\_big} is called with
\(z=x^2\).  For real \(x\), \(\sqrt{z}\) simplifies to \(\operatorname{Abs}(x)\),
but the surrounding Slater/hyperexpand prefactor still contributes a factor of
\(x\).  The expanded primitive therefore contains
\[
  \frac{x\,\operatorname{acosh}(|x|)}{|x|}
  - \frac{i\pi x}{2|x|}.
\]
The imaginary branch constant does not affect the derivative away from zero;
the real multiplier \(x/|x|\) is the part that flips the derivative on
\(x<-1\).  A plausible fix direction supported by the diagnosis is for the
hypergeometric representative expansion to preserve the square-root branch, or
to split by the sign of the real variable rather than only by \(x^2>1\), so that
the outer real intervals reduce to \(\operatorname{acosh}(|x|)\) rather than
\(x\operatorname{acosh}(|x|)/|x|\).  A defensive integration-layer check could
also reject a Meijer indefinite result when differentiating it fails to recover
the integrand under the active assumptions, but the source branch error appears
to be in the hyperexpand representative continuation.

\begin{lstlisting}[style=bugpython]
S.NegativeOne**n*((S.Half - n)*pi/sqrt(z)
                  + I*acosh(sqrt(z))/sqrt(z))
\end{lstlisting}

\subsection{Additional Failing Instances}

The independent verification checks the representative point \(x=-2\) and five
additional negative real inputs on the same smooth branch \(x<-1\).  For each
case, SymPy's returned primitive differentiates to
\(\operatorname{asec}(x)-2/\sqrt{x^2-1}\) rather than to
\(\operatorname{asec}(x)\).  The expected residual is therefore zero for every
row, because all listed points lie away from the branch endpoint and any valid
antiderivative must differentiate back to the original integrand on that
interval.

\begin{center}
\small
\renewcommand{\arraystretch}{1.3}
\begin{tabular}{@{}>{\raggedright\arraybackslash}p{0.44\linewidth}>{\raggedright\arraybackslash}p{0.23\linewidth}>{\raggedright\arraybackslash}p{0.23\linewidth}@{}}
\toprule
\textbf{Input / Expression} & \textbf{SymPy behavior} & \textbf{Expected behavior} \\
\midrule
\(x=-2\), \(\frac{d}{dx}\operatorname{integrate}(\operatorname{asec}(x),x)-\operatorname{asec}(x)\) & residual \(=-2\sqrt{3}/3\) & residual \(=0\) \\
\(x=-3\), same residual expression & residual \(=-\sqrt{2}/2\) & residual \(=0\) \\
\(x=-3/2\), same residual expression & residual \(=-4\sqrt{5}/5\) & residual \(=0\) \\
\(x=-4\), same residual expression & residual \(=-2\sqrt{15}/15\) & residual \(=0\) \\
\(x=-5/2\), same residual expression & residual \(=-4\sqrt{21}/21\) & residual \(=0\) \\
\(x=-10\), same residual expression & residual \(=-2\sqrt{11}/33\) & residual \(=0\) \\
\bottomrule
\end{tabular}
\end{center}

\subsection{Related Issues Assessment}

The bug-finding harness performed a best-effort deduplication check.  This
candidate belongs to a known family.  The \texttt{integrate(asec(x), x)}
capability under test was added by pull request \#19993 (``upgrades to
manualintegrate'', merged 2020-09-02,
\url{https://github.com/sympy/sympy/pull/19993}), which closed issue \#14329
(``Integration of asec and acsc'',
\url{https://github.com/sympy/sympy/issues/14329}) and issue \#19983 (``Sympy
doesn't Integrate when expr contains asec / acsc'',
\url{https://github.com/sympy/sympy/issues/19983}); issue \#14329 already shows
the \texttt{Piecewise}/\(\operatorname{acosh}\) shape whose \texttt{real=True}
variant is the buggy primitive here.  The same wrong-real-branch
Meijer/\(\operatorname{acosh}\) mechanism is reported in the still-open issue
\#10453 (``integrate return wrong answer (inverse trigonometric functions)'',
\url{https://github.com/sympy/sympy/issues/10453}).  The deduplication verdict
is therefore \emph{known family, specifically new instance}: no public issue,
pull request, Stack Overflow post, mailing-list discussion, or release-note
entry was identified for \texttt{integrate(asec(x), x)} producing a
\emph{nonzero} derivative residual on the smooth negative real branch
\(x<-1\).  This exact reproducible integration failure remains individually
actionable and suitable for a dedicated regression test.

\subsection{Proposed SymPy Regression Test}

\medskip

\begin{lstlisting}[style=bugpython]
import pytest

from sympy import S, asec, diff, integrate, simplify, symbols


@pytest.mark.parametrize(
    "value",
    [
        S(-2),
        S(-3),
        -S(3) / 2,
        S(-4),
        -S(5) / 2,
        S(-10),
    ],
)
def test_integrate_asec_antiderivative_on_negative_real_branch(value):
    x = symbols("x", real=True)
    residual = diff(integrate(asec(x), x), x) - asec(x)
    assert simplify(residual.subs(x, value)) == 0
\end{lstlisting}

\end{document}
